% =============================================================================
% Chapter: Dense Pixel-Wise Alignment Methods for Cross-Modal IR-RGB Registration
% =============================================================================
% Intended for inclusion in a graduate-level LaTeX document via \input{} or
% \include{}.  Compile with the main document that defines \documentclass,
% loads packages (amsmath, graphicx, hyperref, …), and includes the
% corresponding .bib file.
% =============================================================================

\chapter{Dense Pixel-Wise Alignment Methods for Cross-Modal IR--RGB Registration}
\label{ch:dense_alignment}

\section{Introduction}
\label{sec:dense_intro}

Fusing infrared (IR) and visible (RGB) imagery is a prerequisite for
numerous applications in surveillance, autonomous driving, and remote
sensing.  A critical first step is \emph{registration}: bringing
corresponding pixels of the two modalities into spatial agreement.
Classical approaches estimate a global geometric transformation---an
affine or projective homography---that maps one image onto the other.
While adequate when the scene is planar or the camera baseline is
negligible, a single global transformation cannot account for parallax,
lens distortion differences between the two sensors, or local
deformations introduced by independent optical paths.

This chapter surveys methods that instead predict a \emph{dense
displacement field}~$\boldsymbol{\phi}: \mathbb{R}^2 \to \mathbb{R}^2$,
assigning every pixel~$(u,v)$ an offset~$(\Delta x, \Delta y)$ such that
the warped pixel location is $(u + \Delta x,\; v + \Delta y)$.  The
resulting deformation is non-parametric---it can model arbitrary
spatially varying misalignment.  We organise the discussion around five
families of architectures: optical-flow networks
(Section~\ref{sec:optflow}), Spatial Transformer Networks
(Section~\ref{sec:stn}), deformable convolutions
(Section~\ref{sec:dcn}), U-Net--based displacement regressors
(Section~\ref{sec:unet_disp}), and the VoxelMorph registration
framework (Section~\ref{sec:voxelmorph}).  For each, we describe the
architecture, the mechanism by which per-pixel displacements are
produced, and the specific advantages and limitations when the two
images to be aligned exhibit the drastic appearance differences inherent
to the IR--RGB pair.


% =====================================================================
\section{Optical Flow--Based Approaches}
\label{sec:optflow}
% =====================================================================

Optical flow estimation seeks to recover the apparent 2D motion field
between two temporally adjacent frames.  Because the output is a dense
$(dx,dy)$ map at every pixel, the same machinery can in principle be
repurposed for spatial registration of an IR--RGB pair.

\subsection{FlowNet and FlowNet\,2.0}

\textbf{FlowNet}~\cite{dosovitskiy2015flownet} was the first work to
cast optical flow as a supervised learning problem solved by a
convolutional neural network (CNN).  Two variants were proposed:
\emph{FlowNetSimple}~(FlowNetS), which concatenates the two input
images along the channel dimension and feeds them through a contracting
encoder followed by an expanding decoder with ``upconvolutional''
layers; and \emph{FlowNetCorr}~(FlowNetC), which processes each image
through a separate feature tower and explicitly computes a correlation
volume between the two feature maps before decoding.  Both produce a
dense flow field at the original resolution via successive upsampling and
refinement steps.  Training was performed on the synthetic \emph{Flying
Chairs} dataset.

\textbf{FlowNet\,2.0}~\cite{ilg2017flownet2} extended the original work
with three key contributions: (i)~a carefully designed training schedule
that progressively increases data complexity, (ii)~a stacked cascade of
FlowNet modules in which each stage warps the second image by the
current flow estimate before refining the residual, and (iii)~a
specialised sub-network for small displacements.  The result was a
halving of endpoint error relative to FlowNet while retaining real-time
throughput.

\subsection{PWC-Net}

PWC-Net~\cite{sun2018pwcnet} embeds three classical optical-flow
principles---\textbf{P}yramid processing, \textbf{W}arping, and
\textbf{C}ost-volume construction---into a compact, end-to-end trainable
CNN.  A learnable feature pyramid is constructed for both images.  At
each pyramid level the second image's features are warped by the
current flow estimate, a partial cost volume is built between the warped
features and the first image's features (limiting the displacement
search range for efficiency), and a lightweight CNN refines the flow.
The coarse-to-fine strategy allows the network to capture large
displacements at coarse levels and fine detail at fine levels.  PWC-Net
is roughly 17$\times$ smaller than FlowNet\,2.0 and achieved
state-of-the-art accuracy on MPI~Sintel and KITTI~2015 at the time of
publication.

\subsection{RAFT}

\textbf{RAFT} (Recurrent All-Pairs Field Transforms)~\cite{teed2020raft}
marked a paradigm shift by replacing coarse-to-fine pyramids with an
iterative, recurrent refinement operating on a single high-resolution
correlation volume.  The architecture comprises three components:

\begin{enumerate}
    \item \textbf{Feature encoder.} A shared-weight residual network
          maps each input image to a feature map at $\nicefrac{1}{8}$
          resolution ($H/8 \times W/8 \times 256$).  A separate
          \emph{context encoder} extracts features from the first image
          only.
    \item \textbf{All-pairs correlation volume.}  The dot product of
          every feature vector in image~1 with every feature vector in
          image~2 produces a 4D $H \times W \times H \times W$ volume.
          The last two dimensions are average-pooled at multiple scales,
          yielding a multi-scale correlation pyramid.
    \item \textbf{Recurrent update operator.}  A convolutional GRU
          iteratively refines the flow estimate.  At each iteration it
          performs a lookup into the correlation pyramid around the
          current correspondence and updates the flow via a learned
          residual.  Because weights are shared across iterations, the
          number of refinement steps can be varied at inference time.
\end{enumerate}

RAFT won the ECCV~2020 Best Paper Award and achieved a 16\% error
reduction on KITTI and 30\% on Sintel (final pass) relative to
prior published results, while using only 4.8\,M parameters.

\subsection{Adaptation to Cross-Modal IR--RGB Registration}

Optical-flow networks are trained on pairs of images with consistent
photometric appearance; the brightness constancy assumption is deeply
embedded in the synthetic training data (Flying Chairs, FlyingThings3D,
Sintel) and in photometric losses.  When the two inputs are an IR and an
RGB frame, this assumption is violated: a hot object may be bright in IR
yet dark in RGB, and vice versa.  Several strategies have been explored
to bridge this domain gap:

\begin{itemize}
    \item \textbf{Cross-modal image translation.}  An auxiliary network
          translates the RGB image into a pseudo-IR image (or vice
          versa), reducing the problem to mono-modal flow estimation.
          Recent work by~\cite{bird2024rgb2ir} generates synthetic IR
          optical-flow datasets by applying an RGB-to-IR style transfer
          to existing flow benchmarks, enabling direct supervision.
    \item \textbf{Feature-level matching.}  Rather than computing
          correlation on raw pixels, one can extract modality-invariant
          features (e.g.\ via a domain-adversarial encoder) and compute
          the cost volume in that shared feature space.
    \item \textbf{Structural similarity losses.}  Replacing the
          photometric $L_1$ or $L_2$ reconstruction loss with
          modality-robust objectives such as mutual information, the
          structural similarity index (SSIM) computed on gradient
          magnitudes, or normalised cross-correlation can guide
          unsupervised fine-tuning on real IR--RGB pairs.
\end{itemize}

\noindent\textbf{Limitations.}  Even with these adaptations, purely
flow-based methods face two challenges: (i)~they are designed for
temporally adjacent frames and may struggle with the large appearance
gap when no temporal cue exists, and (ii)~publicly available
pre-trained weights encode strong inductive biases toward same-modality
matching, making domain transfer non-trivial.


% =====================================================================
\section{Spatial Transformer Networks}
\label{sec:stn}
% =====================================================================

\subsection{The Original STN}

The Spatial Transformer Network (STN), introduced by
Jaderberg~et~al.~\cite{jaderberg2015stn}, provides a differentiable
module that can be inserted into any CNN to perform learned spatial
transformations of feature maps.  The module consists of three
components:

\begin{enumerate}
    \item \textbf{Localisation network.}  A small CNN (or fully
          connected network) that takes the input feature map and
          regresses the parameters~$\boldsymbol{\theta}$ of the desired
          transformation (e.g.\ a $2 \times 3$ affine matrix).
    \item \textbf{Grid generator.}  Given~$\boldsymbol{\theta}$, a
          regular grid of output coordinates is mapped back to input
          coordinates via $\begin{pmatrix} x_i^s \\ y_i^s
          \end{pmatrix} = \mathbf{A}_{\boldsymbol{\theta}}
          \begin{pmatrix} x_i^t \\ y_i^t \\ 1 \end{pmatrix}$, where
          $(x_i^t, y_i^t)$ are target (output) grid coordinates and
          $(x_i^s, y_i^s)$ are the corresponding source (input)
          sampling locations.
    \item \textbf{Differentiable sampler.}  Bilinear (or other
          sub-pixel) interpolation is applied at each $(x_i^s, y_i^s)$
          to produce the output pixel value.  Crucially, the gradient of
          the interpolation kernel with respect to both the input feature
          map and the sampling coordinates is well-defined, enabling
          end-to-end back-propagation through the entire module.
\end{enumerate}

The original paper demonstrated affine and projective transformations,
but the framework is not limited to low-dimensional parametric
families.

\subsection{Extension to Dense Deformation Fields}

The key insight enabling non-parametric warping is that the grid
generator need not be parameterised by a compact~$\boldsymbol{\theta}$.
Instead, the localisation network can directly output a per-pixel
displacement map~$\boldsymbol{\phi} \in \mathbb{R}^{H \times W \times
2}$.  The sampling grid then becomes
\[
    x_i^s = x_i^t + \phi_x(x_i^t, y_i^t), \qquad
    y_i^s = y_i^t + \phi_y(x_i^t, y_i^t),
\]
where $\phi_x, \phi_y$ are the predicted horizontal and vertical
displacements.  The differentiable bilinear sampler remains unchanged;
gradients flow back through the displacements into the localisation
network.

This dense variant underpins most modern learned registration methods.
de~Vos~et~al.~\cite{devos2017dirnet} proposed \textbf{DIRNet}, which
couples a convolutional regressor with a B-spline--parameterised spatial
transformer, outputting a control-point grid that is upsampled to a full
displacement field.  Training is unsupervised, using an image-similarity
loss (e.g.\ normalised cross-correlation) back-propagated through the
sampler.

\subsection{Relevance to IR--RGB Alignment}

For the IR--RGB case, the STN's appeal is that the entire alignment step
is differentiable and can be embedded \emph{inside} a larger
task network (e.g.\ a fusion or detection network).  The downstream task
loss---object detection AP, segmentation IoU, or a fusion quality
metric---propagates gradients back through the sampler into the
displacement predictor, allowing the alignment to be optimised jointly
with the end task.  This is especially valuable when ground-truth
displacement fields are unavailable and only task-level supervision
exists.

\noindent\textbf{Pros:} Fully differentiable, end-to-end trainable with
task loss, trivially extensible from affine to dense warps, lightweight
sampler.

\noindent\textbf{Cons:} The localisation network must be expressive
enough to predict good displacements from cross-modal input; without
guidance from a suitable loss, training may converge to trivial or
identity transformations.  The standard bilinear sampler also introduces
mild blurring, which may degrade high-frequency texture details in the
warped image.


% =====================================================================
\section{Deformable Convolutions}
\label{sec:dcn}
% =====================================================================

\subsection{Deformable ConvNets v1}

Standard convolutions operate on a fixed, regular grid of sampling
locations.  \textbf{Deformable Convolution}~\cite{dai2017dcnv1}
augments each sampling location~$\mathbf{p}_k$ in the kernel with a
learned 2D offset~$\Delta\mathbf{p}_k$, so that the output at
position~$\mathbf{p}_0$ becomes
\[
    y(\mathbf{p}_0) = \sum_{k=1}^{K}
        w_k \cdot x(\mathbf{p}_0 + \mathbf{p}_k + \Delta\mathbf{p}_k),
\]
where $K$ is the number of sampling points in the kernel (e.g.\ $K=9$
for a $3\times 3$ kernel), $w_k$ are the convolution weights, and
$\Delta\mathbf{p}_k$ is predicted by a parallel convolution applied to
the same input feature map.  Because the offsets are typically
fractional, bilinear interpolation is used, making the operation
differentiable.  The effect is that the receptive field of each output
neuron can adapt its shape to the local geometric structure of the
input.

\subsection{Deformable ConvNets v2 (DCNv2)}

\textbf{DCNv2}~\cite{zhu2019dcnv2} introduced two improvements: (i)~the
use of deformable convolutions in many more layers of the backbone
(stages~3--5 of ResNet, totalling 13 layers for ResNet-50), and
(ii)~a \textbf{modulation mechanism} that multiplies each sample by a
learned scalar~$\Delta m_k \in [0,1]$:
\[
    y(\mathbf{p}_0) = \sum_{k=1}^{K}
        w_k \cdot \Delta m_k \cdot
        x(\mathbf{p}_0 + \mathbf{p}_k + \Delta\mathbf{p}_k).
\]
A single convolution layer applied to the input feature map predicts
$3K$ channels: $2K$ for the offsets~$\Delta\mathbf{p}_k$ and $K$ for
the modulation scalars~$\Delta m_k$ (passed through a sigmoid).  The
modulation mechanism allows the network not only to shift its sampling
locations but also to suppress irrelevant or occluded samples entirely.

\subsection{Deformable Convolutions as Implicit Alignment}

In a dual-stream IR--RGB feature extractor, replacing standard
convolutions with deformable convolutions allows each branch to
\emph{implicitly} align its features with the other stream.  Concretely,
suppose features from the IR encoder and the RGB encoder are concatenated
(or cross-attended) at a certain resolution.  A deformable convolution
applied to this joint feature map learns offsets that adaptively sample
from the spatially misaligned IR or RGB features, effectively performing
local geometric correction within the feature space itself, without
producing an explicit displacement field.

This implicit alignment strategy has been successfully deployed in
video super-resolution (e.g.\ EDVR~\cite{wang2019edvr}), where
deformable convolutions align neighbouring frames to the reference
frame at the feature level.  The same principle transfers directly to
IR--RGB alignment: misaligned features from the two modalities are
brought into correspondence through learned offsets before being fused.

\noindent\textbf{Pros:} No explicit warp of the image is needed;
alignment happens in feature space, which is potentially more robust to
appearance differences because high-level features can be
modality-invariant.  Deformable convolutions integrate naturally into
existing backbones (ResNet, HRNet) with minimal architectural change.

\noindent\textbf{Cons:} The alignment is \emph{implicit}---there is no
readily interpretable displacement field to inspect or regularise.
Offset learning can be unstable without careful initialisation and
learning-rate tuning.  Moreover, because offsets are predicted
independently at each spatial location and each layer, global
consistency of the alignment is not guaranteed; the network may
learn offsets that are locally optimal but globally incoherent.


% =====================================================================
\section{U-Net / Encoder--Decoder Architectures for Dense Displacement Prediction}
\label{sec:unet_disp}
% =====================================================================

\subsection{Architecture Overview}

The U-Net architecture~\cite{ronneberger2015unet} is the dominant
backbone for dense prediction tasks in which \emph{every pixel} of the
output must be informed by both local texture and global context.  Its
encoder--decoder structure with skip connections makes it a natural
choice for predicting a two-channel displacement field
$(\Delta x, \Delta y)$ from a pair of input images.

A typical configuration for IR--RGB registration concatenates the IR
image and the RGB image along the channel dimension to form a
$(H \times W \times 4)$ input tensor (one IR channel plus three RGB
channels, or one channel each if the RGB is converted to luminance).
The encoder applies a series of convolutional blocks interleaved with
$2\times$ downsampling (strided convolution or max-pooling), doubling
the feature channels at each level while halving spatial resolution.
The bottleneck captures a compressed, high-level representation of the
joint appearance.  The decoder mirrors the encoder with transposed
convolutions or bilinear upsampling followed by convolution, halving
the channel count and doubling spatial resolution at each level.  A
final $1 \times 1$ convolution projects the decoded features to two
output channels representing $(\Delta x, \Delta y)$.

\subsection{The Role of Skip Connections}

Skip connections are the defining feature that distinguishes U-Net from
a plain encoder--decoder.  At each resolution level, the encoder's
feature map is concatenated with the corresponding decoder feature map
before the next decoder block.  For displacement prediction, these
connections are critical: the encoder's shallow layers contain
high-resolution edge and texture information necessary for sub-pixel
displacement accuracy, while the decoder's deep features encode
semantically informed, large-scale correspondence.  Without skip
connections, fine spatial detail is lost during encoding and cannot be
recovered by the decoder, leading to blurred, overly smooth displacement
fields.

Advanced skip-connection strategies further improve performance.
\textbf{UNet++}~\cite{zhou2018unetpp} introduces nested, dense skip
pathways that progressively bridge the semantic gap between encoder and
decoder features at each resolution.  For displacement regression,
dense skip connections enable multi-scale fusion of deformation cues:
coarse levels capture large translations while fine levels refine
sub-pixel offsets.

\subsection{Loss Functions for Cross-Modal Displacement Regression}

When ground-truth displacement fields are available (e.g.\ from
synthetic warps or from a conventional registration pipeline applied
offline), the network can be trained with a direct regression loss:
\[
    \mathcal{L}_{\text{disp}} = \frac{1}{N}\sum_{i=1}^{N}
        \|\hat{\boldsymbol{\phi}}_i - \boldsymbol{\phi}_i^{*}\|_1,
\]
where $\hat{\boldsymbol{\phi}}_i$ is the predicted displacement at
pixel~$i$ and $\boldsymbol{\phi}_i^{*}$ is the ground-truth
displacement.

In the more common unsupervised setting, the loss combines an
\textbf{image similarity} term and a \textbf{smoothness
regulariser}:
\[
    \mathcal{L} = \mathcal{L}_{\text{sim}}\bigl(
        I_{\text{fixed}},\; I_{\text{moving}} \circ
        \boldsymbol{\phi}\bigr)
        + \lambda\,\mathcal{L}_{\text{smooth}}(\boldsymbol{\phi}).
\]
For cross-modal pairs, conventional photometric losses (MSE, $L_1$) are
inappropriate because pixel intensities are not comparable across
modalities.  Suitable alternatives include:
\begin{itemize}
    \item \textbf{Normalised mutual information (NMI)}, which measures
          statistical dependence between intensity distributions
          regardless of their functional relationship;
    \item \textbf{Modality-independent neighbourhood descriptor
          (MIND)}~\cite{heinrich2012mind}, which computes a
          self-similarity context around each pixel, yielding a
          modality-invariant representation that can be compared with
          SSD;
    \item \textbf{Gradient-based SSIM}, computed on the magnitude of
          image gradients rather than raw intensities, exploiting the
          observation that edges tend to co-localise across IR and RGB.
\end{itemize}
The smoothness term is typically a diffusion regulariser
$\mathcal{L}_{\text{smooth}} = \sum_i \|\nabla \boldsymbol{\phi}_i\|^2$,
penalising large spatial gradients in the displacement field to
encourage locally smooth, physically plausible deformations.

\subsection{Practical Considerations}

\begin{itemize}
    \item \textbf{Multi-scale prediction.}  Predicting displacement
          fields at multiple decoder levels and composing them
          (coarse-to-fine) improves convergence and captures both large
          and small deformations~\cite{zhao2019recursive}.
    \item \textbf{Diffeomorphic constraints.}  Integrating a velocity
          field via scaling-and-squaring ensures the deformation is a
          diffeomorphism (smooth, invertible), preventing
          folding~\cite{dalca2019unsupervised}.
    \item \textbf{Data augmentation.}  Given only a few thousand aligned
          pairs, applying random thin-plate-spline warps to one modality
          creates a virtually unlimited supply of training pairs with
          known ground-truth displacements.
\end{itemize}

\noindent\textbf{Pros:} Flexible, well-understood architecture; skip
connections preserve fine spatial detail; directly outputs an
interpretable displacement field; large ecosystem of pre-trained
encoders.

\noindent\textbf{Cons:} Requires careful choice of cross-modal
similarity loss; with only a few thousand pairs, overfitting is a
concern without aggressive augmentation; the concatenation-based input
strategy treats the two modalities symmetrically, which may not be
optimal when their information content differs substantially.


% =====================================================================
\section{VoxelMorph and Medical Image Registration Parallels}
\label{sec:voxelmorph}
% =====================================================================

\subsection{The VoxelMorph Framework}

VoxelMorph~\cite{balakrishnan2019voxelmorph} was developed for
deformable registration of 3D medical volumes (e.g.\ brain MRI) but
its architectural principles are directly transferable to 2D IR--RGB
registration.  The framework concatenates a \emph{fixed} image and a
\emph{moving} image along the channel dimension, passes the result
through a 3D U-Net encoder--decoder, and outputs a dense displacement
field~$\boldsymbol{\phi}$ at the input resolution.  A differentiable
spatial transformer (extending Jaderberg~et~al.'s 2D sampler to 3D with
trilinear interpolation) applies~$\boldsymbol{\phi}$ to the moving
image.

Training is unsupervised.  The loss function has two terms:
\[
    \mathcal{L} = \mathcal{L}_{\text{sim}}(I_{\text{fixed}},\;
        I_{\text{moving}} \circ \boldsymbol{\phi})
        + \lambda \sum_{\mathbf{p}} \|\nabla\boldsymbol{\phi}(\mathbf{p})\|^2,
\]
where $\mathcal{L}_{\text{sim}}$ is mean squared error or normalised
cross-correlation, and the second term is a diffusion regulariser that
encourages smooth deformations.  An optional \emph{auxiliary
segmentation loss} can be included when anatomical labels are available,
using the Dice coefficient between the warped moving-image segmentation
and the fixed-image segmentation.

\subsection{Diffeomorphic Variant}

A diffeomorphic extension~\cite{dalca2019unsupervised} replaces the
direct displacement output with a stationary velocity
field~$\mathbf{v}$.  The displacement field is obtained by integrating
$\mathbf{v}$ through time via the scaling-and-squaring algorithm:
\[
    \boldsymbol{\phi} = \exp(\mathbf{v})
        \approx \underbrace{(\mathrm{id} + \mathbf{v}/2^T)
        \circ \cdots \circ (\mathrm{id} + \mathbf{v}/2^T)}_{T \text{
        compositions}},
\]
where $T$ is the number of integration steps (typically 7).  The
resulting deformation is guaranteed to be a diffeomorphism (smooth and
invertible), which prevents the biologically implausible ``folding'' of
tissue that can occur with unconstrained displacement fields.  The
integration is implemented as a sequence of spatial transformer layers
and is fully differentiable.

\subsection{Application to 2D IR--RGB Registration}

Adapting VoxelMorph from 3D medical volumes to 2D IR--RGB pairs is
straightforward:

\begin{enumerate}
    \item Replace all 3D convolutions with 2D convolutions and the
          trilinear sampler with a bilinear sampler.
    \item Set the input to be a two-channel (or four-channel) tensor
          formed by concatenating the single-channel IR image with the
          RGB image (or its luminance channel).
    \item Replace the similarity loss with a cross-modal metric
          (NMI, MIND, or gradient-domain NCC) as discussed in
          Section~\ref{sec:unet_disp}.
    \item Optionally enforce diffeomorphism via velocity-field
          integration if large deformations are expected.
\end{enumerate}

The VoxelMorph codebase\footnote{\url{https://github.com/voxelmorph/voxelmorph}}
provides modular implementations of the spatial transformer, the
integration layer, and various loss functions in PyTorch and TensorFlow,
making it an accessible starting point.

\subsection{Comparison with Direct U-Net Regression}

VoxelMorph can be viewed as a specific instantiation of the U-Net
displacement regressor discussed in Section~\ref{sec:unet_disp}, with
two notable additions: (i)~the explicit spatial transformer layer that
warps the moving image and feeds the warped result into the loss, and
(ii)~the optional diffeomorphic integration layer.  The conceptual
contribution of VoxelMorph is less about architectural novelty and more
about establishing a principled, reproducible framework for
learning-based registration with well-defined loss functions,
regularisation strategies, and evaluation protocols.

In the IR--RGB context, the VoxelMorph paradigm offers a disciplined
approach to balancing alignment accuracy (via the similarity loss) and
deformation regularity (via the smoothness penalty).  The
hyperparameter~$\lambda$ controlling this trade-off can be tuned on a
held-out validation set using application-specific metrics (e.g.\ edge
alignment error, landmark transfer accuracy, or downstream fusion
quality).

\noindent\textbf{Pros:} Well-tested, open-source framework;
diffeomorphic option guarantees topology preservation; modular design
separates the displacement predictor from the warping and loss
computation; strong community and extensive ablation studies in the
literature.

\noindent\textbf{Cons:} Originally designed for mono-modal (MRI--MRI)
registration; cross-modal application requires replacing the similarity
loss, which may degrade convergence.  The assumption of a single,
globally smooth deformation field may be insufficient if the IR and RGB
sensors have independent, spatially varying distortions.


% =====================================================================
\section{Discussion and Practical Recommendations}
\label{sec:discussion}
% =====================================================================

Table~\ref{tab:method_comparison} summarises the five method families.
For a practitioner with a few thousand aligned IR--RGB pairs, we offer
the following recommendations:

\begin{enumerate}
    \item \textbf{Start with a U-Net displacement regressor
          (Section~\ref{sec:unet_disp}) or VoxelMorph
          (Section~\ref{sec:voxelmorph}).}  These architectures are
          simple, well-understood, and directly output interpretable
          displacement fields.  Using synthetic random warps for data
          augmentation can effectively multiply the training set.
    \item \textbf{Use a cross-modal similarity loss.}  MIND descriptors
          or gradient-domain NCC are robust to the appearance gap.
          Avoid raw MSE between IR and RGB intensities.
    \item \textbf{Consider optical-flow pre-training.}  Initialise the
          encoder with weights from a RAFT or PWC-Net model pre-trained
          on standard flow benchmarks, then fine-tune on IR--RGB data.
          The feature encoder learns generally useful correspondence
          primitives that transfer across domains.
    \item \textbf{Use deformable convolutions for implicit alignment
          in a fusion backbone.}  If the end goal is a downstream task
          (detection, segmentation), integrating DCNv2 layers into the
          fusion network allows the alignment to be optimised jointly
          with the task, avoiding the need for a separate registration
          step.
    \item \textbf{Regularise aggressively.}  With limited data, the
          diffusion smoothness penalty and, where appropriate,
          diffeomorphic integration are essential to prevent
          degenerate displacement fields.
\end{enumerate}

\begin{table}[t]
\centering
\caption{Comparison of dense pixel-wise alignment method families for
IR--RGB registration.}
\label{tab:method_comparison}
\small
\begin{tabular}{p{3.0cm} p{2.8cm} p{2.8cm} p{2.5cm} p{2.5cm}}
\hline
\textbf{Method} & \textbf{Displacement type} & \textbf{Cross-modal readiness} & \textbf{Key strength} & \textbf{Key weakness} \\
\hline
Optical flow (RAFT, PWC-Net) & Explicit dense flow & Low (needs adaptation) & Mature, accurate & Brightness constancy assumed \\
STN (dense variant) & Explicit dense warp & Medium & End-to-end with task & May collapse to identity \\
Deformable Conv (DCNv2) & Implicit per-layer offsets & Medium--High & No explicit warp needed & Not interpretable \\
U-Net displacement & Explicit dense field & High (with proper loss) & Simple, interpretable & Needs cross-modal loss \\
VoxelMorph & Explicit dense field & High (with proper loss) & Diffeomorphic option & Designed for mono-modal \\
\hline
\end{tabular}
\end{table}


% =====================================================================
\section{Conclusion}
\label{sec:dense_conclusion}
% =====================================================================

Dense pixel-wise alignment of IR and RGB imagery is a challenging
problem that sits at the intersection of optical flow estimation,
learned image registration, and cross-modal feature matching.  The five
families of methods surveyed in this chapter offer complementary
strengths.  Optical-flow networks provide mature, high-accuracy dense
correspondence but require domain adaptation for cross-modal inputs.
Spatial Transformer Networks supply the differentiable warping primitive
upon which all learned registration methods build.  Deformable
convolutions offer an implicit, feature-level alignment that integrates
seamlessly into task-driven pipelines.  U-Net--based displacement
regressors and the VoxelMorph framework provide principled,
interpretable dense deformation prediction with well-studied
regularisation strategies.  For the scenario of a few thousand aligned
IR--RGB pairs, a U-Net or VoxelMorph-style regressor trained with a
modality-robust similarity loss and augmented with synthetic warps
represents a practical and effective starting point, with deformable
convolutions offering a complementary path when the alignment is to be
optimised jointly with a downstream task.


% =============================================================================
% Bibliography entries (to be placed in a .bib file)
% =============================================================================
% The following \cite{} keys are used throughout this chapter.  The
% corresponding BibTeX entries are provided in the companion file
% references_dense_alignment.bib.
% =============================================================================
